\documentclass[11pt]{article}
\usepackage[a4paper,margin=1in]{geometry}
\usepackage[T1]{fontenc}
\usepackage{lmodern}
\usepackage{amsmath,amssymb,amsthm,mathtools,mathrsfs}
\usepackage{array,booktabs,microtype,xurl,placeins}
\usepackage{hyperref}
\hypersetup{colorlinks=true,linkcolor=black,citecolor=black,urlcolor=blue,
  pdftitle={A Trace-SOS Determinantal Inequality for Partitions of the Identity},
  pdfauthor={Zixuan Ye and GPT-6 Astra}}
\allowdisplaybreaks[1]
\setlength{\emergencystretch}{1em}
\newtheorem{theorem}{Theorem}[section]
\newtheorem{lemma}[theorem]{Lemma}
\newtheorem{proposition}[theorem]{Proposition}
\newcommand{\tr}{\operatorname{tr}}
\newcommand{\Rea}{\operatorname{Re}}
\newcommand{\HS}{\mathrm{HS}}
\newcommand{\Id}{I_n}
\newcommand{\cK}{\mathcal K}
\newcommand{\cB}{\mathcal B}
\newcommand{\gH}{\mathscr H}
\newcommand{\gQ}{\mathscr Q}
\newcommand{\ip}[2]{\langle #1,#2\rangle}
\newcommand{\nrm}[1]{\lVert #1\rVert}
\newcommand{\word}[1]{\mathtt{#1}}

\title{\textbf{A Trace--SOS Determinantal Inequality\\
for Partitions of the Identity}}
\author{Zixuan Ye \and GPT-6 Astra}
\date{}

\begin{document}
\maketitle
\begin{abstract}
Let \(A_1,\ldots,A_m\) be positive semidefinite Hermitian \(n\times n\)
matrices whose sum is the identity. We prove
\[
 \det\!\left(\Id+m(m-1)\sum_{i=1}^m A_i^3\right)
 \ge (2m-1)^n\det\!\left(\sum_{i=1}^m A_i^2\right).
\]

The argument extends Cubesteak's trace--sum-of-squares approach to the
three-matrix case. Its main ingredient is an explicit family of Gram
matrices, valid for every \(m\ge3\), together with an analytic
decomposition proving their positive semidefiniteness. In fact, we obtain
a stronger quadratic trace inequality and the corresponding inverse-trace
bound. The constant is sharp. For positive definite matrices and
\(m\ge3\), equality holds precisely when all matrices equal \(\Id/m\);
for \(m=1,2\), the determinant inequality is an identity. A single
optional Python certificate verifies the polynomial trace identity
using exact rational arithmetic and an indeterminate parameter.
\end{abstract}

\section{Introduction and main result}

Cubesteak proved the following inequality on Mathematics Stack Exchange
\cite{cubesteak}: if \(A,B,C\) are positive definite Hermitian matrices
with \(A+B+C=\Id\), then
\[
 \det\!\bigl(\Id+6(A^3+B^3+C^3)\bigr)
 \ge 5^n\det(A^2+B^2+C^2).
\]

The proof compares determinants through an inverse-trace inequality,
removes the inverse by a quadratic variational identity, and establishes
positivity by a noncommutative sum-of-squares decomposition.
We retain this proof framework and its notation \(S,T,\Phi\).
The new ingredient is a uniform Gram family that treats an arbitrary
number of summands.

Throughout, \(m,n\) are positive integers, \(\Id\) is the identity in
\(M_n(\mathbb C)\), and \(A_i=A_i^*\succeq0\) satisfy
\[
 \sum_{i=1}^m A_i=\Id.
\]

We use the Hilbert--Schmidt inner product and norm
\[
 \ip{U}{V}_{\HS}=\tr(U^*V),\qquad
 \nrm{U}_{\HS}^2=\tr(U^*U).
\]

Set
\begin{equation}\label{eq:ST}
 c_m=m(m-1),\qquad d_m=2m-1,\qquad
 S=\sum_i A_i^2,\qquad R=\sum_i A_i^3,\qquad
 T=\frac{\Id+c_mR}{d_m}.
\end{equation}

These definitions are valid on the positive semidefinite boundary.
Indeed,
\begin{equation}\label{eq:positive-ST}
 S-\frac1m\Id=\sum_i\left(A_i-\frac1m\Id\right)^2\succeq0,
 \qquad T\succeq\frac1{d_m}\Id,
\end{equation}
so \(S\) and \(T\) are always positive definite.

\begin{theorem}\label{thm:main}
For the matrices above,
\begin{equation}\label{eq:trace-bound}
 \tr(T^{-1}S)\le n
\end{equation}
and
\begin{equation}\label{eq:det-bound}
 \det\!\left(\Id+m(m-1)\sum_i A_i^3\right)
 \ge(2m-1)^n\det\!\left(\sum_i A_i^2\right).
\end{equation}
The constant \((2m-1)^n\) is sharp for each \(m,n\).
If every \(A_i\) is positive definite and \(m\ge3\), equality in
\eqref{eq:det-bound} holds if and only if
\(A_1=\cdots=A_m=\Id/m\).
For \(m=1,2\), both \eqref{eq:trace-bound} and
\eqref{eq:det-bound} are equalities for every admissible tuple.
\end{theorem}

The positive definiteness assumption in the equality classification
matters. For example, when \(m\ge2\), the positive semidefinite tuple
\[
 A_1=0,\qquad A_2=\cdots=A_m=\frac1{m-1}\Id
\]
also gives equality, since \(S=T=\Id/(m-1)\).
We do not classify all equality cases on the boundary.

\section{The trace--variational reduction}

For a Hermitian matrix \(X\), define
\begin{equation}\label{eq:KB}
 \cK(X)=\sum_i A_iXA_i,\qquad
 \cB(X)=\tr(XTXS)-\nrm{\cK(X)}_{\HS}^2
\end{equation}
and
\begin{equation}\label{eq:Phi}
 \Phi(X)=\tr(\Id-2XS+XTXS).
\end{equation}

The expression \(\tr(XTXS)\) is real and nonnegative, since it equals
\(\nrm{T^{1/2}XS^{1/2}}_{\HS}^2\).

\begin{lemma}\label{lem:reduction}
For every Hermitian \(X\),
\begin{align}
 \Phi(X)&=\nrm{\cK(X)-\Id}_{\HS}^2+\cB(X),
 \label{eq:Phi-K}\\
 \Phi(X)&=n-\tr(T^{-1}S)
 +\nrm{T^{1/2}(X-T^{-1})S^{1/2}}_{\HS}^2.
 \label{eq:Phi-square}
\end{align}
Consequently, if \(\cB(X)\ge0\) for all Hermitian \(X\), then
\eqref{eq:trace-bound} and \eqref{eq:det-bound} hold.
\end{lemma}

\begin{proof}
Since \(\cK(X)\) is Hermitian and
\(\tr\cK(X)=\tr(XS)\), expansion of the squared norm proves
\eqref{eq:Phi-K}. Put \(Y=X-T^{-1}\). Expansion of \(XTX\) gives
\[
 XTX=YTY+2Y+T^{-1},
\]
which proves \eqref{eq:Phi-square}. In particular,
\(\Phi(T^{-1})=n-\tr(T^{-1}S)\). Nonnegativity in
\eqref{eq:Phi-K} therefore implies \eqref{eq:trace-bound}.

For the determinant conclusion, the matrix
\(M=T^{-1/2}ST^{-1/2}\) is positive definite. The arithmetic--geometric
mean inequality for its eigenvalues gives
\begin{equation}\label{eq:AMGM}
 \left(\frac{\det S}{\det T}\right)^{1/n}
 =(\det M)^{1/n}
 \le\frac{\tr M}{n}
 =\frac{\tr(T^{-1}S)}{n}\le1.
\end{equation}

Using \eqref{eq:ST} yields \eqref{eq:det-bound}.
\end{proof}

We shall prove the following stronger assertion.
\begin{theorem}\label{thm:quadratic}
For every admissible tuple \(A_1,\ldots,A_m\) and every Hermitian \(X\),
\begin{equation}\label{eq:quadratic}
 \cB(X)\ge0.
\end{equation}
\end{theorem}

The rest of the proof constructs a Gram representation of
\(d_m\cB(X)\) and proves positivity of its coefficient matrix
analytically.

\section{A uniform trace--Gram identity}

Fix \(m\ge3\) and write \(k=m-1\ge2\). For each \(i\), let
\(J_i=\{1,\ldots,m\}\setminus\{i\}\). For \(j,l\in J_i\), define
\begin{align}
 E_{ij}&=(A_j-A_i)XA_j,\qquad
 F_{ij}=A_jX(A_j-A_i)=E_{ij}^*, \label{eq:EF}\\
 Z_{i;jl}&=(A_j-A_i)X(A_l-A_i)\qquad(j\ne l).
 \label{eq:Z}
\end{align}

Let \(w_i\) be the column of all \(E_{ij}\), all \(F_{ij}\), and all
\(Z_{i;jl}\) with ordered pairs \(j\ne l\), in any fixed order.
It has \(2k+k(k-1)=k(k+1)\) components.
The corresponding abstract component labels are \(E_j,F_j,Z_{jl}\).

Define a real symmetric matrix \(\gH_k\) on these labels by
Table~\ref{tab:gram}; entries in symmetric positions are equal.
The table is unchanged by a common relabeling of the \(k\) indices.
It defines a new parameterized family, distinct from the fixed Gram
matrix in Cubesteak's three-matrix proof.

\begin{table}[htbp]
\centering
\caption{The entries of \(\gH_k\). Every \(Z\)-label has distinct internal
indices. The conditions in each block are exhaustive.}
\label{tab:gram}
\renewcommand{\arraystretch}{1.18}
\begin{tabular}{lll}
\toprule
Row and column & Condition & Entry\\
\midrule
\(E_j,E_j\) & & \(k(k+5)/2\)\\
\(E_j,E_l\) & \(j\ne l\) & \((k-3)/2\)\\
\(F_j,F_j\) & & \(k(k+5)/2\)\\
\(F_j,F_l\) & \(j\ne l\) & \((5k+1)/2\)\\
\(E_j,F_j\) & & \(-(k+1)(k+2)/2\)\\
\(E_j,F_l\) & \(j\ne l\) & \(-(k+1)\)\\
\(E_j,Z_{lh}\) & \(j=l\) & \(k/2\)\\
\(E_j,Z_{lh}\) & \(j\ne l\) & \(-1/2\)\\
\(F_j,Z_{lh}\) & \(j=l\) & \(-(k-1)/2\)\\
\(F_j,Z_{lh}\) & \(j=h\) & \(0\)\\
\(F_j,Z_{lh}\) & \(j\notin\{l,h\}\) & \(1/2\)\\
\(Z_{jl},Z_{jl}\) & & \(k-1\)\\
\(Z_{jl},Z_{ph}\) & distinct labels; \(j=p\) or \(l=h\) & \(-1/2\)\\
\(Z_{jl},Z_{ph}\) & distinct labels; \(j\ne p,\ l\ne h\) & \(0\)\\
\bottomrule
\end{tabular}
\end{table}

For a Hermitian matrix \(A\) and a column \(w\) of matrices indexed by
these labels, put
\begin{equation}\label{eq:gram-functional}
 \gQ_k(A;w)=\sum_{a,b}(\gH_k)_{ab}\tr(w_a^*Aw_b).
\end{equation}

The sum is real, by Hermitian symmetry.

\begin{proposition}\label{prop:identity}
For every Hermitian \(X\),
\begin{equation}\label{eq:universal-gram}
 d_m\cB(X)=\sum_{i=1}^m\gQ_k(A_i;w_i).
\end{equation}
Equivalently,
\begin{equation}\label{eq:expanded-target}
 \sum_i\gQ_k(A_i;w_i)
 =\tr(X^2S)+m(m-1)\tr(XRXS)
   -(2m-1)\nrm{\cK(X)}_{\HS}^2.
\end{equation}
This algebraic identity requires only Hermitian symmetry and
\(\sum_iA_i=\Id\), not positivity of the \(A_i\).
\end{proposition}

\begin{proof}
Expand \eqref{eq:gram-functional} using
\eqref{eq:EF}--\eqref{eq:Z} and Table~\ref{tab:gram}.
Appendix~\ref{app:identity} gives the full coefficient table, together
with the finite index-summation rules used to obtain it.
The resulting three nonzero coefficients are \(1,m(m-1)\), and
\(1-2m\), on the respective trace words in
\eqref{eq:expanded-target}. This proves that equation.
Finally, \(d_mT=\Id+c_mR\) gives \eqref{eq:universal-gram}.
\end{proof}

\section{Analytic positivity of the Gram family}

We prove positivity for vectors in an arbitrary complex Hilbert space
\(V\), with inner product linear in its second argument.
For a vector of components
\[
 \boldsymbol{\xi}=(\ell_j,r_j,z_{jl})_{1\le j,l\le k,\ j\ne l},
\]
where the \(\ell\), \(r\), and \(z\) entries have the respective
\(E\), \(F\), and \(Z\) labels, write
\[
 Q_k(\boldsymbol{\xi})
 =\sum_{a,b}(\gH_k)_{ab}\ip{\xi_a}{\xi_b}.
\]

All norms and inner products in this section are those of \(V\).

\subsection{An orthogonal decomposition of the ordered-pair components}

Define
\[
 \ell_0=\frac1k\sum_j\ell_j,\qquad
 r_0=\frac1k\sum_jr_j,\qquad
 z_0=\frac1{k(k-1)}\sum_{j\ne l}z_{jl},
 \qquad
 \ell_j^\circ=\ell_j-\ell_0,\quad r_j^\circ=r_j-r_0.
\]

Let
\[
 u_j=\sum_{l\ne j}z_{jl}-(k-1)z_0,\qquad
 v_j=\sum_{l\ne j}z_{lj}-(k-1)z_0.
\]

Both \(\sum_ju_j\) and \(\sum_jv_j\) vanish. Put
\begin{equation}\label{eq:alpha-beta}
 \alpha_j=
 \begin{cases}
 (u_j+v_j)/(2(k-2)),&k>2,\\
 0,&k=2,
 \end{cases}
 \qquad
 \beta_j=\frac{u_j-v_j}{2k},
\end{equation}
and
\begin{equation}\label{eq:residual}
 \omega_{jl}=z_{jl}-z_0-\alpha_j-\alpha_l-\beta_j+\beta_l
 \qquad(j\ne l).
\end{equation}

For \(k=2\), one has \(u_j+v_j=0\), so all formulas below still hold.

We have \(\sum_j\alpha_j=\sum_j\beta_j=0\), and every row sum and
column sum of the residual array is zero. In particular,
\begin{align}
 a_j:=\sum_{l\ne j}z_{jl}
  &=(k-1)z_0+(k-2)\alpha_j+k\beta_j, \label{eq:rows}\\
 b_j:=\sum_{l\ne j}z_{lj}
  &=(k-1)z_0+(k-2)\alpha_j-k\beta_j. \label{eq:columns}
\end{align}

The four arrays in
\[
 z_{jl}=z_0+(\alpha_j+\alpha_l)+(\beta_j-\beta_l)+\omega_{jl}
\]
are mutually orthogonal in the direct-sum Hilbert space indexed by
ordered pairs \(j\ne l\).
Indeed, the zero sums give orthogonality to the constant array.
The \(\alpha\)-array is symmetric and the \(\beta\)-array is
antisymmetric under \(j\leftrightarrow l\).
Orthogonality to the residual array follows from its zero row and
column sums. Expanding the two remaining squared norms yields
\begin{equation}\label{eq:pair-norm}
 \sum_{j\ne l}\nrm{z_{jl}}^2
 =k(k-1)\nrm{z_0}^2
 +2(k-2)\sum_j\nrm{\alpha_j}^2
 +2k\sum_j\nrm{\beta_j}^2
 +\sum_{j\ne l}\nrm{\omega_{jl}}^2.
\end{equation}

\subsection{Completion of squares}

For a real symmetric \(2\times2\) matrix \(N\), abbreviate
\[
 \ip{(a,b)}{N(a,b)}
 =N_{11}\nrm a^2+2N_{12}\Rea\ip a b+N_{22}\nrm b^2.
\]
Define
\begin{align}
 N_0&=\frac14
 \begin{pmatrix}
 4k^2+k+7&-(6k^2+5k+1)\\
 -(6k^2+5k+1)&12k^2+k-1
 \end{pmatrix},\label{eq:N0}\\
 N_1&=\frac1{k+2}
 \begin{pmatrix}
 (k+1)(2k+3)&-(k+1)^2\\
 -(k+1)^2&k^2+k-1
 \end{pmatrix}.\label{eq:N1}
\end{align}

\begin{lemma}\label{lem:Gram-positive}
For every integer \(k\ge2\),
\begin{align}
 Q_k(\boldsymbol{\xi})
 ={}&k(k-1)\nrm{z_0+\tfrac12(\ell_0-r_0)}^2
       +k\ip{(\ell_0,r_0)}{N_0(\ell_0,r_0)}
       \notag\\
 &+(k-2)(k+2)\sum_j
    \nrm{\alpha_j+\frac{k+1}{2(k+2)}
                    (\ell_j^\circ-r_j^\circ)}^2
       \notag\\
 &+k^2\sum_j
    \nrm{\beta_j+
             \frac{(k+1)\ell_j^\circ-(k-1)r_j^\circ}{2k}}^2
       \notag\\
 &+\sum_j\ip{(\ell_j^\circ,r_j^\circ)}
                   {N_1(\ell_j^\circ,r_j^\circ)}
       +k\sum_{j\ne l}\nrm{\omega_{jl}}^2.
 \label{eq:completed-Gram}
\end{align}
Consequently, \(\gH_k\succeq0\) for \(k\ge2\), and
\(\gH_k\succ0\) for \(k\ge3\).
For \(k=2\), the equality \(Q_2(\boldsymbol{\xi})=0\) implies
\begin{equation}\label{eq:k2-kernel}
 \ell_j^\circ=r_j^\circ=0\quad(j=1,2),
 \qquad 5\ell_0=7r_0.
\end{equation}
\end{lemma}

\begin{proof}
For a transparent verification of \eqref{eq:completed-Gram}, put
\(\ell_\Sigma=\sum_j\ell_j\), \(r_\Sigma=\sum_jr_j\), and
\(z_\Sigma=\sum_{j\ne l}z_{jl}\).
Reading off the six blocks of Table~\ref{tab:gram} gives
\begin{align}
 Q_k(\boldsymbol{\xi})={}&
 \frac{(k+1)(k+3)}2\sum_j\nrm{\ell_j}^2
   +\frac{k-3}2\nrm{\ell_\Sigma}^2
 \notag\\
 &+\frac{(k-1)(k+1)}2\sum_j\nrm{r_j}^2
   +\frac{5k+1}2\nrm{r_\Sigma}^2
 \notag\\
 &-k(k+1)\Rea\sum_j\ip{\ell_j}{r_j}
   -2(k+1)\Rea\ip{\ell_\Sigma}{r_\Sigma}
 \notag\\
 &+k\sum_{j\ne l}\nrm{z_{jl}}^2
   -\frac12\sum_j\bigl(\nrm{a_j}^2+\nrm{b_j}^2\bigr)
 \notag\\
 &+\Rea\left((k+1)\sum_j\ip{\ell_j}{a_j}
                    -\ip{\ell_\Sigma}{z_\Sigma}\right)
 \notag\\
 &+\Rea\left(\ip{r_\Sigma}{z_\Sigma}
              -k\sum_j\ip{r_j}{a_j}-\sum_j\ip{r_j}{b_j}\right).
 \label{eq:uncompleted}
\end{align}

For instance, the \(ZZ\)-block has diagonal \(k-1\) and off-diagonal
entry \(-1/2\) for equal rows or equal columns. Its contribution is
exactly the fourth line of \eqref{eq:uncompleted}.
The two subsequent lines are the \(EZ\)- and \(FZ\)-contributions.

Substitute \eqref{eq:rows}--\eqref{eq:pair-norm} and
\(\ell_j=\ell_0+\ell_j^\circ\), \(r_j=r_0+r_j^\circ\)
into \eqref{eq:uncompleted}. Terms between constant and zero-sum
components vanish. The coefficients of
\(\nrm{z_0}^2,\sum_j\nrm{\alpha_j}^2,
\sum_j\nrm{\beta_j}^2\) are, respectively,
\(k(k-1),(k-2)(k+2),k^2\).
Completing these three squares leaves the constant block \(kN_0\),
the zero-sum block \(N_1\), and
\(k\sum_{j\ne l}\nrm{\omega_{jl}}^2\), exactly as in
\eqref{eq:completed-Gram}.

Both upper-left entries in \eqref{eq:N0}--\eqref{eq:N1} are positive,
and direct determinant expansion gives
\begin{equation}\label{eq:two-determinants}
 \det N_0=\frac{(k-2)(k-1)^2(3k+1)}4,\qquad
 \det N_1=\frac{(k+1)(k^2-2)}{k+2}.
\end{equation}
Thus \(N_0\succeq0\) and \(N_1\succ0\) for \(k\ge2\);
both are positive definite for \(k\ge3\). All terms in
\eqref{eq:completed-Gram} are nonnegative, proving
\(\gH_k\succeq0\).

If \(k\ge3\) and \(Q_k(\boldsymbol{\xi})=0\), the \(N_0\)- and
\(N_1\)-terms force
\(\ell_0=r_0=\ell_j^\circ=r_j^\circ=0\).
The other squares then force \(z_0=\alpha_j=\beta_j=\omega_{jl}=0\).
Hence every component is zero, proving strict positivity.
For \(k=2\),
\[
 N_0=\frac14\begin{pmatrix}25&-35\\-35&49\end{pmatrix}.
\]

The \(N_1\)-terms force \(\ell_j^\circ=r_j^\circ=0\);
the \(N_0\)-term forces \(5\ell_0=7r_0\).
This proves the necessary conditions \eqref{eq:k2-kernel}.
\end{proof}

\section{Proofs of the inequalities and the equality statement}

\begin{proof}[Proof of Theorem~\ref{thm:quadratic}]
Suppose first that \(m\ge3\).
Apply Lemma~\ref{lem:Gram-positive} in the Hilbert space
\(M_n(\mathbb C)\) with its Hilbert--Schmidt inner product, to the
components \(A_i^{1/2}(w_i)_a\). Then
\[
 \gQ_k(A_i;w_i)
 =\sum_{a,b}(\gH_k)_{ab}
       \ip{A_i^{1/2}(w_i)_a}{A_i^{1/2}(w_i)_b}_{\HS}\ge0.
\]

Proposition~\ref{prop:identity} and \(d_m>0\) prove \(\cB(X)\ge0\).

For \(m=1\), \(A_1=S=T=\Id\), so \(\cB(X)=0\).
For \(m=2\), write \(A_2=\Id-A_1\). Direct expansion gives
\[
 \Id+2\bigl(A_1^3+(\Id-A_1)^3\bigr)
 =3\bigl(A_1^2+(\Id-A_1)^2\bigr),
\]
and hence \(T=S\). The two matrices \(A_i\) are simultaneously
diagonalizable. In a common orthonormal eigenbasis write their
diagonal entries as \(a_{ip}\), and put
\[
 s_p=\sum_{i=1}^2a_{ip}^2,\qquad
 h_{pq}=\sum_{i=1}^2a_{ip}a_{iq}.
\]

For Hermitian \(X=(x_{pq})\),
\[
 \tr(XSXS)=\sum_{p,q}s_ps_q|x_{pq}|^2,\qquad
 \nrm{\cK(X)}_{\HS}^2=\sum_{p,q}h_{pq}^2|x_{pq}|^2.
\]

Cauchy--Schwarz gives \(h_{pq}^2\le s_ps_q\).
Therefore \(\cB(X)\ge0\) also for \(m=2\).
\end{proof}

\begin{proof}[Proof of Theorem~\ref{thm:main}]
The two inequalities follow from Theorem~\ref{thm:quadratic} and
Lemma~\ref{lem:reduction}. The preceding proof also shows \(S=T\)
for \(m=1,2\), giving the asserted identities in these cases.
For any \(m\), the choice \(A_i=\Id/m\) gives
\[
 S=T=\frac1m\Id,
\]
so equality holds and the constant is sharp.

It remains to prove uniqueness of equality in the positive definite
domain when \(m\ge3\). Assume \(\det S=\det T\).
Both inequalities in \eqref{eq:AMGM} must be equalities. In particular,
\(\tr(T^{-1}S)=n\). Set \(X=T^{-1}\succ0\).
Equations \eqref{eq:Phi-K} and \eqref{eq:universal-gram} give
\[
 0=\Phi(X)=\nrm{\cK(X)-\Id}_{\HS}^2
       +\frac1{d_m}\sum_i\gQ_k(A_i;w_i).
\]

Every summand is nonnegative; hence
\(\gQ_k(A_i;w_i)=0\) for each \(i\).

If \(m\ge4\), then \(k\ge3\) and \(\gH_k\succ0\).
It follows that \(A_i^{1/2}E_{ij}=0\) for every \(j\ne i\).
Since \(A_i\) is positive definite, \(E_{ij}=0\); thus
\[
 (A_j-A_i)XA_j=0.
\]

Both \(X\) and \(A_j\) are invertible, so \(A_j=A_i\).
All matrices are therefore \(\Id/m\).

If \(m=3\), then \(k=2\). Apply \eqref{eq:k2-kernel} to
\(\ell_j=A_i^{1/2}E_{ij}\) and
\(r_j=A_i^{1/2}F_{ij}\), with the two elements of \(J_i\)
identified with \(1,2\).
Canceling the invertible factor \(A_i^{1/2}\), the two
\(E_{ij}\)'s have a common value \(E_i^0\), the two
\(F_{ij}\)'s have a common value \(F_i^0\), and
\[
 5E_i^0=7F_i^0.
\]

By \eqref{eq:EF}, \(F_i^0=(E_i^0)^*\). Taking adjoints also gives
\(5F_i^0=7E_i^0\), whence \(E_i^0=F_i^0=0\).
The same invertibility argument yields \(A_j=A_i\), completing the
equality classification.
\end{proof}

\appendix
\section{Exact expansion of the trace--Gram identity}
\label{app:identity}

This appendix provides the algebraic details of
Proposition~\ref{prop:identity}. All reductions use only real trace
identities, finite sums, and \(\sum_iA_i=\Id\).

\subsection{Trace words and the complete coefficient table}

For a word \(W\) in the symbols \(A,B,X\), define \(\Theta(W)\) as
follows. Replace each distinct symbol among \(A,B\) occurring in \(W\)
by an independently summed matrix \(A_i\), leave \(X\) fixed, and take
the real part of the trace. For example,
\[
 \Theta(\word{AAAXBBX})
   =\sum_{i,j}\Rea\tr(A_i^3XA_j^2X)=\tr(XRXS),
 \qquad
 \Theta(\word{AAXX})=\tr(X^2S).
\]

Different summation labels are independent and may have equal values.
Cyclic rotation, reversal, and consistent renaming of the summed
symbols leave \(\Theta\) unchanged. Reversal is valid because every
letter represents a Hermitian matrix and the real trace is taken.
Also,
\[
 \Theta(\word{ABXBAX})
 =\sum_{i,j}\Rea\tr(A_iA_jXA_jA_iX)
 =\nrm{\cK(X)}_{\HS}^2.
\]

Expand \(\sum_i\gQ_k(A_i;w_i)\) and separate its six contributions
according to component types. The columns \(EF,EZ,FZ\) each include
both orders of the corresponding cross terms. Table~\ref{tab:expansion}
lists every nonzero word coefficient after setting \(k=m-1\).
A blank entry means zero; all words not listed have zero coefficient.
The entries are polynomials in \(m\), not values at selected integers.

\begin{table}[!htbp]
\centering
\caption{Complete trace-word expansion of the six Gram blocks.}
\label{tab:expansion}
\begingroup
\setlength{\tabcolsep}{3.5pt}
\renewcommand{\arraystretch}{1.65}
\small
\begin{tabular}{lrrrrrr}
\toprule
Word & \(EE\) & \(EF\) & \(EZ\) & \(FF\) & \(FZ\) & \(ZZ\)\\
\midrule
\(\word{AAAXBBX}\)&
 \(\frac{m^2}2+m\)&&\(-m\)&&&\(\frac{m^2}2-m\)\\
\(\word{AAAXX}\)&
 \(\frac m2-2\)&&\(1\)&&&\(1-\frac m2\)\\
\(\word{AABXABX}\)&&
 \(m-m^2\)&\(m^2-m\)&&&\\
\(\word{AABXBBX}\)&
 \(-m^2-2m\)&\(m^2-m\)&\(2m\)&&&\(m\)\\
\(\word{AABXBX}\)&
 \(4-m\)&\(2m\)&\(-2-m\)&&&\(-2\)\\
\(\word{AAXAAX}\)&
 \(\frac{m^2}2+m\)&\(m-m^2\)&\(-m\)&
 \(\frac{m^2}2-m\)&\(m\)&\(-m\)\\
\(\word{AAXAX}\)&&
 \(-2m\)&\(1\)&\(-2+\frac{5m}2\)&\(-2\)&\(3-\frac m2\)\\
\(\word{AAXBABX}\)&&&&
 \(\frac{m^2}2-m\)&\(2m-m^2\)&\(\frac{m^2}2-m\)\\
\(\word{AAXX}\)&&&\(1\)&&&\\
\(\word{ABAXABX}\)&&
 \(m^2-m\)&\(m-m^2\)&\(2m-m^2\)&\(m^2-3m\)&\(m\)\\
\(\word{ABAXBX}\)&&&&&\(m-1\)&\(1-m\)\\
\(\word{ABXABX}\)&&&&&\(1\)&\(-1\)\\
\(\word{ABXBAX}\)&
 \(\frac m2-2\)&&\(m\)&\(2-\frac{5m}2\)&\(2-m\)&\(-1\)\\
\(\word{AXAX}\)&&&\(-1\)&&&\(1\)\\
\bottomrule
\end{tabular}
\endgroup
\end{table}

Adding across each row leaves only
\[
 m(m-1)\Theta(\word{AAAXBBX})
 +\Theta(\word{AAXX})
 +(1-2m)\Theta(\word{ABXBAX}),
\]
which is \eqref{eq:expanded-target}.

\subsection{Finite summation rules underlying the table}

For completeness, we specify an exact procedure for obtaining all
entries of Table~\ref{tab:expansion}, including the treatment of
distinct indices. Before expansion, a term has one distinguished index
\(i\) and at most four other index positions. Separate these terms
according to the equality pattern of the latter positions, keeping
\(i\) distinct from all of them. Each pattern is counted once and is
summed over all assignments of pairwise distinct indices to its
distinct labels. The Gram coefficient is determined by
Table~\ref{tab:gram} for that pattern.

The following finite inclusion--exclusion identity changes such
distinct-index sums into unrestricted sums. For any function \(f\) of
\(r\) indices in \(\{1,\ldots,m\}\),
\begin{equation}\label{eq:partition-inversion}
 \sum_{\substack{i_1,\ldots,i_r\\\text{pairwise distinct}}}
 f(i_1,\ldots,i_r)
 =
 \sum_{\pi\in\Pi_r}
 \left[\prod_{C\in\pi}(-1)^{|C|-1}(|C|-1)!\right]
 \sum_{\substack{i_1,\ldots,i_r\\
       i_a=i_b\ \text{if }a,b\in C\in\pi}}
 f(i_1,\ldots,i_r).
\end{equation}

Here \(\Pi_r\) denotes the set of set partitions of \(\{1,\ldots,r\}\);
indices belonging to different blocks in the inner sum remain
independent and are allowed to coincide.

To prove \eqref{eq:partition-inversion}, fix an assignment of all the
indices and let \(\rho\) be its equality partition. Its coefficient on
the right is the sum of the displayed partition weights over
\(\pi\) refining \(\rho\). This factors over the blocks of \(\rho\).
For a block of size \(s\), the factor is the sum of the signs of all
permutations of \(s\) objects: a partition specifies their cycle
supports, and a block of size \(t\) has \((t-1)!\) cyclic orderings,
each contributing sign \((-1)^{t-1}\).
The sign sum is \(1\) for \(s=1\) and \(0\) for \(s\ge2\), since
composition with a fixed transposition pairs opposite signs.
Thus the total coefficient is \(1\) exactly for pairwise distinct
indices and \(0\) otherwise, as required.

After applying \eqref{eq:partition-inversion}, expand each
\(E,F,Z\) into its two or four monomials. Every resulting trace word
initially has five \(A\)-factors and two \(X\)-factors, because each
entry of \(w_i\) has two \(A\)-factors and one \(X\)-factor and
\eqref{eq:gram-functional} inserts one additional \(A_i\).
An unrestricted summed label appearing exactly once can be removed:
summing its single matrix factor replaces it by
\(\sum_iA_i=\Id\). A summed label absent from the word contributes
a scalar factor \(m\). Consequently at most two matrix labels remain,
as each remaining label occurs at least twice among at most five
matrix factors.

Finally identify words under cyclic rotation, reversal, and renaming
of summed labels, choosing the lexicographically least representative
after renaming labels by order of first occurrence. Expanding with
these rules and collecting the six component types yields precisely
Table~\ref{tab:expansion}. All operations involve at most five index
labels, so \eqref{eq:partition-inversion} uses only \(r\le5\).
This argument is valid also when \(r>m\), in which case the
distinct-index sum is empty and the right side cancels to zero.
Thus the table and the Gram identity hold for every integer
\(m\ge3\) without an interpolation or limiting argument.

\FloatBarrier
\section{The accompanying exact certificate}\label{app:certificate}

The sole accompanying program is
\begin{center}
 \texttt{stage6\_round2\_identity\_certificate.py}.
\end{center}

It independently checks \eqref{eq:universal-gram} by implementing the
finite expansion rules of Appendix~\ref{app:identity}.
Polynomial coefficients are stored exactly using rational numbers;
the parameter \(m\) is indeterminate. Four auxiliary labels suffice
to enumerate every equality pattern of a pair of component labels;
the distinguished index gives at most five labels in
\eqref{eq:partition-inversion}. The program collects the coefficients
of the normalized trace words and checks that their sum is exactly
\[
 \Theta(\word{AAXX})+
 (m^2-m)\Theta(\word{AAAXBBX})+
 (1-2m)\Theta(\word{ABXBAX}).
\]

It uses only the Python standard library. Running
\begin{verbatim}
python3 stage6_round2_identity_certificate.py
\end{verbatim}
produces
\begin{verbatim}
PASS: symbolic identity for indeterminate m
\end{verbatim}

The certificate performs no matrix sampling, floating-point
optimization, or extrapolation from finitely many values of \(m\).
It is an optional verification of the algebraic identity.
Positivity of the Gram family, the trace reduction, and the equality
classification are proved analytically in the paper.

\begin{thebibliography}{9}
\bibitem{cubesteak}
Cubesteak,
\emph{Answer concerning the determinant inequality
\(\det(\Id+6(A^3+B^3+C^3))\ge5^n\det(A^2+B^2+C^2)\)},
Mathematics Stack Exchange, answer no.~5147902.
\url{https://math.stackexchange.com/a/5147902}.
\end{thebibliography}
\end{document}

